% Source for /assets/pdf/cv.pdf. Rebuild with:
%   tectonic _cv_source/cv.tex --outdir assets/pdf

% Document class and font size
\documentclass[a4paper,9pt]{extarticle}

% Packages
\usepackage[utf8]{inputenc}
\usepackage{geometry}
\geometry{letterpaper, margin=0.75in}
\usepackage{titlesec}
\usepackage{enumitem}
\usepackage[hidelinks]{hyperref}
\usepackage{tabularx}
\usepackage{fancyhdr}

% Formatting
\setlist{noitemsep}
\titleformat{\section}{\large\bfseries}{\thesection}{1em}{}[\titlerule]
\titlespacing*{\section}{0pt}{\baselineskip}{\baselineskip}

% Begin document
\begin{document}

\pagestyle{fancy}
\renewcommand{\headrulewidth}{0pt}
\fancyhead{}
\fancyhead[L]{\textit{Prathamesh Devadiga}}
\fancyhead[R]{\textit{September 2026}}
\thispagestyle{empty}

% Header
\begin{flushleft}
\textbf{\LARGE Prathamesh Devadiga}\\[2pt]
CS PhD Student, JASPR Lab, Dartmouth College
\\ \href{mailto:prathamesh.p.devadiga.gr@dartmouth.edu}{prathamesh.p.devadiga.gr@dartmouth.edu} | \href{https://www.linkedin.com/in/prathamesh-devadiga}{linkedin.com/in/prathamesh-devadiga} | \href{https://prathameshdevadiga.vercel.app}{prathameshdevadiga.vercel.app}
\end{flushleft}

% Research Interests
\section*{RESEARCH INTERESTS}
\noindent
Machine learning security and privacy, memorization and training-data extraction in LLMs, machine unlearning and model editing auditing, adversarial robustness, LLM jailbreaks and defenses, production AI systems

% Education Section
\section*{EDUCATION}

\noindent
\textbf{Dartmouth College}, Hanover, NH, USA \hfill Aug 2026 -- Present\\
Doctor of Philosophy (PhD) in Computer Science\\
Advisor: Prof. Shawn Shan, JASPR Lab (Joint AI Security Privacy Research)\\
Research Interests: Trustworthy Machine Learning, AI Security and Privacy, Large Language Models, Memorization and Generalization Theory, Adversarial Machine Learning, Responsible AI Systems

\vspace{0.5em}

\noindent
\textbf{PES University}, Bangalore, India \hfill Aug 2022 -- May 2026\\
Bachelor of Technology in Computer Science Engineering \hfill Cumulative GPA: 8.98/10.00\\
Relevant Coursework: Machine Learning, Deep Learning, Generative AI, Linear Algebra and its Applications, Statistics for Data Science, Engineering Mathematics I \& II, Data Structures and Algorithms, Operating Systems, Computer Networks, Distributed Systems

% Research Experience
\section*{RESEARCH EXPERIENCE}

\noindent
\textbf{Dartmouth College} \hfill Remote\\
\textit{Research Assistant (ML Security \& LLM Privacy) — Advisor: Prof. Shawn Shan} \hfill July 2025 -- July 2026
\begin{itemize}
    \item Co-developed \textbf{CHASE (Correlated Hierarchical Adaptive Search for Extraction)}, a Gaussian Process bandit framework for query-efficient memorization extraction from instruction-tuned large language models
    \item Discovered that the prompt space for memorization extraction exhibits a \textbf{low-rank, kernel-correlated structure}, enabling effective prompt optimization with as few as \textbf{100 API queries} per target passage
    \item Evaluated CHASE across \textbf{9 instruction-tuned LLMs from 7 model families}, including production APIs, demonstrating up to \textbf{2.9$\times$} higher extraction success than random search while maintaining over \textbf{10$\times$} gains under seven defense strategies
    \item Showed that ordinary academic and continuation prompts consistently outperform adversarial jailbreaks for memorization extraction, highlighting limitations of existing prompt-level safety defenses
\end{itemize}

\noindent
\textbf{Adhāra AI Labs} \hfill Remote\\
\textit{Founder \& Lead Researcher} \hfill Feb 2025 -- Present
\begin{itemize}
    \item Founded independent research laboratory focused on Generative AI, Large Language Models (LLMs), Retrieval-Augmented Generation (RAG), and deep learning systems
    \item Leading research initiatives from concept to production, ensuring results translate into scalable, real-world applications
    \item Research outputs include PyraFuseNet (ICIAI), SLMs as Compiler Experts (NeurIPS Workshop), RegimeNAS, MorphNAS, and a low-latency jailbreak prevention framework
\end{itemize}

\noindent
\textbf{Lossfunk} \hfill Bangalore, India\\
\textit{AI Research Intern} \hfill Sep 2025 -- June 2026
\begin{itemize}
    \item Built an Indic LLM pipeline for extremely low-resource languages (Tulu: $\sim$0.001\% of typical training data)
    \item Developed a framework using hard negative constraints to reduce catastrophic language leakage from 80\% to 5\%, demonstrating that explicit prohibitions outperform positive instructions
    \item Paper accepted at the EACL 2026 Workshop on Language Models for Low-Resource Languages (LoResLM): ``Making Large Language Models Speak Tulu: Structured Prompting for an Extremely Low-Resource Language''
\end{itemize}

\noindent
\textbf{Nokia} \hfill Bengaluru, India\\
\textit{Intent-Based Network Management System (3GPP SON)} \hfill Aug 2025 -- Oct 2025
\begin{itemize}
    \item Designed an intent-based self-organizing network (SON) system enabling operators to express high-level network goals in natural language, automatically translated into 3GPP-compliant intents
    \item Built an LLM-powered intent classification pipeline using Retrieval-Augmented Generation (RAG) to dynamically map natural language inputs to structured intent schemas
    \item Architected an intent orchestration layer supporting synchronous and asynchronous execution, conflict detection, duplicate intent suppression, and intent expiry management
    \item Developed a REST-based workflow integrating an interactive web frontend, intent classifier, orchestrator, and Auto Pilot (AP) APIs for automated network configuration
\end{itemize}

\noindent
\textbf{Google Summer of Code 2025} \hfill Remote\\
\textit{University of California Santa Cruz} \hfill May 2025 -- Aug 2025
\begin{itemize}
    \item Selected for highly competitive global program (acceptance rate less than 10\%)
    \item Built billion-scale vector embedding benchmarks (768, 1024, 2048 dimensions) from open-source codebases
    \item Developed infrastructure for evaluating Approximate Nearest Neighbor (ANN) algorithms at scale under realistic workloads
\end{itemize}

\noindent
\textbf{Indian Institute of Technology, Indore} \hfill Remote\\
\textit{Undergraduate Research Intern (Deep Learning) — Advisor: Prof. Nagendra Kumar} \hfill Dec 2023 -- Aug 2024
\begin{itemize}
    \item Designed KASPER framework achieving 99.5\% accuracy in PDF malware detection with adversarial robustness
    \item Developed custom malware injection pipeline for realistic payload embedding and comprehensive model evaluation
    \item Maintained 99.5\% detection accuracy against FGSM and PGD adversarial attacks
    \item Led benchmarking and explainability analysis using Kolmogorov-Arnold Networks (KANs)
    \item Contributions formed core technical implementation for paper accepted at Applied Soft Computing (Q1 Journal)
\end{itemize}

\section*{PUBLICATIONS}
\noindent
\textbf{Conference and Workshop Papers}
\begin{itemize}

\item \textbf{Devadiga, P.} (2026). \textit{Deletion Scars: Inferring Deleted Artist Styles from Paired Diffusion Models}. \textit{ECCV 2026 Workshop on Unlearning and Model Editing (U\&ME)}.

\item \textbf{Devadiga, P.}, Lakshman, A. (2026). \textit{Resistant to Lawyers, Defeated by Disagreement: Evaluation Blindspots in Legal Language Models}. \textit{ICML 2026 Workshop on AI for Law (AI4Law)}.

\item \textbf{Devadiga, P.}, et al. (2026). \textit{Making Large Language Models Speak Tulu: Structured Prompting for an Extremely Low-Resource Language}. \textit{EACL 2026 Workshop on Language Models for Low-Resource Languages (LoResLM)}.

\item \textbf{Devadiga, P.}, et al. (2026). \textit{GUARDIAN: Multi-Agent Defense System for Large Language Model Security}. \textit{International Conference on Artificial Intelligence and Information (ICIAI), Waseda University, Japan}.

\item \textbf{Devadiga, P.}, et al. (2025). \textit{SLMs as Compiler Experts: Auto-Parallelization for Heterogeneous Systems}. \textit{NeurIPS 2025 Workshop on Machine Learning for Systems}.

\item \textbf{Devadiga, P.}, et al. (2025). \textit{PyraFuseNet: Dual-Path Network for Resource-Constrained Vision}. \textit{International Conference on Artificial Intelligence and Information (ICIAI), Nanyang Technological University, Singapore}.

\end{itemize}

\noindent\textbf{Preprints}
\begin{itemize}

\item \textbf{Devadiga, P.}, et al. (2025). \textit{SAMVAD: Multi-Agent Framework for Modeling Judicial Reasoning in Indian Jurisprudence}. arXiv:2509.03793.

\item \textbf{Devadiga, P.}, et al. (2025). \textit{RegimeNAS: Regime-Aware Neural Architecture Search for Financial Trading}. arXiv:2508.11338.

\item \textbf{Devadiga, P.}, et al. (2025). \textit{MorphNAS: Differentiable Neural Architecture Search for Multilingual Named Entity Recognition}. arXiv:2508.15836.

\end{itemize}

\noindent\textbf{Media Coverage}
\begin{itemize}

\item Featured in \textit{The Times of India} for research on enabling large language models to generate Tulu through structured prompting (2026).

\item Featured on the \textit{AIM Network} podcast to discuss AI for low-resource languages and structured prompting techniques for Tulu (2026).

\end{itemize}

% Projects Section
\section*{KEY PROJECTS}
\noindent
\textbf{Cerebrum: Production LLM Training and Serving System} \hfill 2025\\
\textit{Independent Research Project}
\begin{itemize}
    \item Built end-to-end LLM system featuring distributed FSDP/DDP training, Flash Attention v2, and Mixture-of-Experts support
    \item Developed \textbf{Mixture-of-Refusals (MoR)}: defense mechanism achieving 2-3x speedup on safe queries while maintaining identical safety guarantees through conditional safety routing
    \item Implemented vLLM-based inference engine with quantization, speculative decoding, and prefix caching; deployed with Docker, Kubernetes, and Prometheus
\end{itemize}

\noindent
\textbf{Arcane ML: Distributed Training Framework} \hfill 2024\\
\textit{Open Source Project}
\begin{itemize}
    \item Built framework for distributed ML training across SSH clusters, Modal Cloud GPUs, and local setups, with PyTorch DDP support and a unified CLI
\end{itemize}

\section*{TEACHING EXPERIENCE}
\noindent
\textbf{PES University} \hfill Bangalore, India\\
\textit{Teaching Assistant — Machine Learning \& Deep Learning} \hfill Sep 2025 -- May 2026
\begin{itemize}
    \item Taught and mentored a cohort of 75+ students across Machine Learning and Deep Learning courses
    \item Prepared lecture slides, delivered tutorial sessions, and conducted weekly hands-on laboratory sessions
    \item Designed and graded assignments, lab exercises, and hackathons to reinforce practical understanding
\end{itemize}

\noindent
\textbf{PESU I/O} \hfill Bangalore, India\\
\textit{Subject Matter Expert — Deep Learning} \hfill Oct 2024 -- Nov 2024
\begin{itemize}
    \item Co-instructed a 30-hour course on Deep Learning from Scratch for a cohort of 40 students
    \item Delivered 20 hours of in-person instruction covering neural networks, CNNs, RNNs, LSTMs, GANs, and Explainable AI
\end{itemize}

\section*{AWARDS \& HONORS}
\noindent
\textbf{Amazon AI-ML Scholar} \hfill 2024\\
Selected from top 1,000 AI/ML students in India

\medskip
\noindent
\textbf{Oxford Machine Learning Summer School} \hfill 2024\\
Selected participant, University of Oxford

\medskip
\noindent
\textbf{Cohere AI Summer School} \hfill 2024\\
Selected participant in program on large language models and NLP

\medskip
\noindent
\textbf{Winner, Cisco ThingQbator Hackathon 6.0} \hfill Nov 2023\\
National hackathon win among 1,000+ competing teams in IoT and emerging technologies

\medskip
\noindent
\textbf{Merit Scholarships (6x)} \hfill 2022--2025\\
MRD Scholarship (2x), CN Rao Scholarship (2x), Distinction Scholarship (2x) for top 5--20\% academic performance, PES University

\section*{LEADERSHIP \& SERVICE}
\noindent
\textbf{Community Mentor \& Technical Speaker}, Independent \hfill 2023 -- Present
\begin{itemize}
    \item Mentored 50+ teams in hackathons, research projects, and WiDS Datathon competitions
    \item Delivered technical talks on the DSPy framework and LLM fine-tuning with LoRA at FOSS United and GDSC events
\end{itemize}

\noindent
\textbf{Head of Technology}, Entrepreneurship Club, PES University \hfill 2023 -- 2024
\begin{itemize}
    \item Led 10+ member technical team; developed club website and managed technical operations for 100+ members
\end{itemize}

\noindent
\textbf{Webmaster}, IEEE Robotics and Automation Society, PES University \hfill 2023 -- 2024

\noindent
\textbf{Core Member}, Google Developer Student Clubs \& Hacker Space, PES University \hfill 2022 -- 2024

% Skills Section
\section*{TECHNICAL SKILLS}
\begin{itemize}
    \item \textbf{Programming Languages:} Python, Go, Julia, SQL, Bash
    \item \textbf{ML/AI Frameworks:} PyTorch, TensorFlow, Hugging Face, vLLM, LangChain, DSPy, MLflow, FastAI
    \item \textbf{Systems \& Infrastructure:} Docker, Kubernetes, Apache Spark, Apache Kafka, Hadoop, Modal, distributed training (FSDP/DDP)
\end{itemize}

\end{document}
